% !TEX root = ../../../main.tex
\section{Multiplicativity and singularity}
\label{sec:linear-algebra-i:multiplicativity-and-singularity}

% Sources: IMG_9265--IMG_9266; October 8 PDF, pp. 5--8 and 18--19.
\subsection{Composition of endomorphisms}

\begin{theorem}[Determinant of a composition]
  \label{thm:linear-algebra-i:determinant-composition}
  For linear maps \(L_1,L_2:V\to V\),
  \[
    \det(L_2\circ L_1)=(\det L_2)(\det L_1).
  \]
\end{theorem}

\begin{proof}
  Choose a basis \((e_1,\ldots,e_n)\) and a top-degree form \(D\)
  normalized by \(D(e_1,\ldots,e_n)=1\). Applying
  \autoref{eq:linear-algebra-i:determinant-pullback} first to \(L_2\)
  and then to \(L_1\) gives
  \begin{align*}
    \det(L_2\circ L_1)
      &=D(L_2(L_1(e_1)),\ldots,L_2(L_1(e_n)))\\
      &=(\det L_2)D(L_1(e_1),\ldots,L_1(e_n))\\
      &=(\det L_2)(\det L_1)D(e_1,\ldots,e_n)\\
      &=(\det L_2)(\det L_1).
  \end{align*}
\end{proof}
% PDF p.18 drops det(L1) in one intermediate line. The preceding
% pullback identity and final line give the complete calculation above.

\subsection{Products of matrices}

\begin{theorem}[Determinant of a matrix product]
  \label{thm:linear-algebra-i:determinant-product}
  If \(A,B\) are \(n\)-by-\(n\) matrices, then
  \[
    \det(AB)=(\det A)(\det B).
  \]
\end{theorem}

\begin{proof}
  For fixed \(A\), define
  \[
    D_A(v_1,\ldots,v_n)=D(Av_1,\ldots,Av_n),
  \]
  where \(D\) is the normalized determinant function on columns.
  Since multiplication by \(A\) is linear, \(D_A\) is antisymmetric
  and multilinear. The one-dimensionality of the space of top-degree
  forms gives
  \[
    D_A=D_A(e_1,\ldots,e_n)D=(\det A)D.
  \]
  Write the columns of \(B\) as \(C_1,\ldots,C_n\).
  The corresponding columns of \(AB\) are \(AC_1,\ldots,AC_n\). Hence
  \[
    \det(AB)=D_A(C_1,\ldots,C_n)=(\det A)(\det B).
  \]
\end{proof}

\subsection{Invertibility and zero determinant}

\begin{theorem}[Determinant and invertibility]
  \label{thm:linear-algebra-i:determinant-invertibility}
  An endomorphism \(L:V\to V\) is singular if and only if \(\det L=0\).
  If \(L\) is invertible, then
  \[
    \det(L^{-1})=\frac{1}{\det L}.
  \]
\end{theorem}

\begin{proof}
  Suppose \(L\) is singular. Choose \(v\ne0\) in \(\ker L\), set
  \(e_1=v\), and extend this vector to a basis \((e_1,\ldots,e_n)\).
  Normalize \(D\) on this basis. Then
  \[
    \det L=D(L(e_1),\ldots,L(e_n))
      =D(0,L(e_2),\ldots,L(e_n))=0.
  \]
  If \(L\) is invertible, multiplicativity and \(\det I=1\) give
  \[
    (\det L^{-1})(\det L)=\det(L^{-1}\circ L)=1.
  \]
  Thus \(\det L\ne0\), and division gives the reciprocal formula.
  These two implications prove the equivalence.
\end{proof}

For a square matrix \(M\), this becomes
\[
  \det M=0\quad\Longleftrightarrow\quad
  \operatorname{rank}M<n\quad\Longleftrightarrow\quad
  M\text{ is singular}.
\]
There is also a direct calculation using the adapted bases of
\autoref{eq:linear-algebra-i:kernel-first-normal-form}.
Writing \(r=\operatorname{rank}M\), place the domain basis vectors
with independent nonzero images first. Suitable invertible basis matrices
\(E,F\) then satisfy
\[
  ME=FJ_r,\qquad M=FJ_rE^{-1},
\]
where \(J_r\) is \(n\)-by-\(n\), with its first \(r\) diagonal entries
equal to \(1\) and all other entries zero. Its determinant is zero for
\(r<n\) and is one for \(r=n\). Since \(\det F\) and \(\det E^{-1}\) are nonzero,
the product formula gives the same criterion.
